\section{Architecture: multimodalidad nativa y agentes visuales}

\subsection{De KiMi VL a K2.5}
KiMi mantuvo en sus primeras versiones el modelo de lenguaje puro separado del modelo VL, pero a partir de K2.5 dejó de hacer esa distinción; el docente usó «multimodalidad nativa» para describir esa arquitectura que, desde el día 0, funde los tokens visuales, las llamadas a herramientas y el texto. K2.5 ya no es «generar primero el lenguaje y luego meter un fotograma», sino que hace que la visión, el texto y las herramientas aprendan juntos dentro del mismo transformer, tratando la attention como un canal multi-agent.

\begin{importantbox}{La evolución de la multimodalidad nativa}
\begin{itemize}
    \item Versiones anteriores: K2, KiMi VL y Qwen3-VL separaban el modelo de lenguaje del modelo de visión, y dependían de late fusion para lograr la correspondencia entre imagen y texto.
    \item K2.5: desde el día 0 junta imágenes, fotogramas de video, líneas de comandos y llamadas a herramientas; la early fusion hace que cada capa tenga una señal modal conjunta.
    \item Esto significa que la «visión» no es una capacidad añadida, sino una dimensión central que recorre toda la attention del transformer; el prompt incluso llega a especificar que «los tokens visuales aparecen antes que los tokens de lenguaje».
\end{itemize}
\end{importantbox}

\subsection{Tokenization, Memory y Early Fusion}
K2.5 lleva más allá el diseño de la tokenization y la memory para que sean conscientes de la multimodalidad: los fotogramas visuales se dividen en patch tokens, se les añade una posición temporal y se mezclan con los tokens de lenguaje al entrar en la cache. El docente mencionó que la early fusion no consiste en arrojar todos los tokens juntos a la attention, sino en usar primero un `modality router` que hace un soft gate, para asegurar que la información visual se remuestree en cada capa y no quede opacada porque el peso del texto sea demasiado grande.

\begin{knowledgebox}{Las tres capas de la estrategia de early fusion}
\begin{itemize}
    \item Usar modality routing para equilibrar los tokens visuales y de lenguaje, manteniendo en la misma layer una proporción visual-texto de 2:1.
    \item Diseñar la temporal memory cache con una estructura de `frame bundle + tool trace`, para facilitar que una later evaluation pueda reproducir un agent loop.
    \item Diseñar la attention mask como «vision first», de modo que los tokens visuales puedan entrar primero al pipeline de puntuación de GRM antes de decidir si se continúa con el siguiente paso.
\end{itemize}
\end{knowledgebox}

\subsection{La autoorganización de la Visual Agency Intelligence}
El posicionamiento central de K2.5 es la «Visual Agency Intelligence»: no solo entiende el contenido visual, sino que además puede iniciar un agent loop para ejecutar tareas, llamar herramientas e incluso controlar la GUI. El docente mencionó `Activation the Vision-Agentic Capability` (esa línea que aparece en la guía), señalando que el modelo, durante la inferencia hacia adelante, debe ver al mismo tiempo los tokens de visión, lenguaje y herramientas, y a la vez que percibe, va proponiendo el siguiente paso de acción.

\begin{figure}[H]
\centering
\includegraphics[width=0.78\textwidth,trim=0 120 0 120,clip]{cover.jpg}
\caption{En la imagen de portada aparecen a la vez el docente, el diagrama de la arquitectura Visual Agentic y la línea de tiempo de K2.5, lo que subraya la presentación conjunta de la visión y el agent loop.\protect\footnotemark}
\end{figure}
\footnotetext{Intervalo de tiempo en el video: 00:32:10--00:32:48, la imagen muestra al docente presentando el agent swarm superpuesto con la multimodalidad nativa.}

\subsection{Agent loop y GUI agents}
En el plano de la architecture, la visual agency no se queda solo en la attention, sino que toma el control de la acción real mediante el agent loop: el modelo planea con VRL (Visual Reinforcement Loop), primero genera la intent, luego llama a la herramienta y por último regresa a Allcastrater. Los GUI agents son una extensión de este plano: se dedican específicamente a simular la interacción con la interfaz, para asegurar que la architecture pueda controlar botones, controles deslizantes, tablas e incluso pestañas del navegador mediante un prompt estandarizado.

\begin{knowledgebox}{Los límites de tarea de los GUI agents}
\begin{itemize}
    \item Los GUI agents se encargan de la percepción y la operación de la interfaz: convierten un screenshot en eventos de mouse y teclado.
    \item El agent de herramientas se encarga de las llamadas a la API: extrae el state de la API response y lo escribe en el tool trace.
    \item El agent de texto se encarga de las salidas explicativas, asegurando que el GRM pueda puntuar de manera unificada a partir de las entradas multimodales.
\end{itemize}
\end{knowledgebox}

\subsection{Interpretabilidad y enrutamiento de señales}
La multimodalidad nativa hace que las attention layers observen a la vez imagen, texto y tokens de herramientas, lo que facilita seguir cada paso que da el modelo. Esta vía también sirve a la explainability: por ejemplo, en los registros de entrenamiento se marcan el «vision tracker» y el «tool calling track», encadenando la señal de reward, la intención de la operación y la puntuación del GRM, lo que nos ayuda a localizar problemas con rapidez al depurar el modelo. El docente señaló en particular que, con una timeline compuesta por GRM scores + tool trace, se puede reproducir en 10 segundos una acción generada aleatoriamente por un agent.

\subsection{Resumen de la sección}
La arquitectura de K2.5 termina por unificar por completo lo que antes eran el lenguaje y la visión separados; la early fusion y el agent-centric design hacen de la «Visual Agency» la forma de trabajo predeterminada del modelo, tal como subraya la figura: la visión, el lenguaje y las herramientas danzan de manera sincronizada.
